\begin{theorem}[Real-time work law for processes]\label{ar:thm:law}
Let $\Pp$ be such a process. In real time, the optimal work per episode at vanishing error is $\kB T\ln2\,\Kb(\Pp)$. Moreover, for every fixed $\delta>0$, any implementation whose controller and reservoirs start in Gibbs states and that draws at most $N\kB T\ln2\,(\Kb(\Pp)-\delta)$ of deterministic work over $N$ episodes, meaning that its battery ends in a definite energy level with probability tending to one, has error tending to one.
\end{theorem}

As for channels, bounding only the mean work gives no such converse: a device that runs the optimal implementation with probability $\lambda$ and does nothing otherwise spends $\lambda$ times the rate and errs with probability about $1-\lambda$. Nor can the Gibbs initial state of the controller simply be dropped (Section~\ref{ar:sec:converse}).

\begin{theorem}[What is free]\label{ar:thm:free}
For a process $\Pp$ as above the following are equivalent:
\begin{enumerate}[label=(\roman*)]
\item $\Kb(\Pp)=0$;
\item driven by fresh Gibbs inputs, implementations of $\Pp$ can release outputs nearly in their product Gibbs state while leaving their bath nearly uncorrelated with all of them: for every $\delta>0$ some implementation within $\delta$ of $\Pp$ has total Gibbs defect at most $\delta$ (Section~\ref{ar:sec:free});
\item $\Pp$ is a limit of causally ordered thermal implementations: one heat bath in its Gibbs state, retained throughout, and a sequence of unitaries that each conserve energy exactly, with every output released on time;
\item the real-time work rate of $\Pp$ is zero.
\end{enumerate}
Each condition forces covariance under the common time translation, but not at each slot separately.
\end{theorem}

Milz et al.\ take Gibbs-preserving channels as the free channels and ask which higher-order transformations preserve them~\cite{Milz26}. Requiring this completely, they find exactly the causally ordered combs that are thermal in their sense: inserting the Gibbs state into the first slot returns the Gibbs state, uncorrelated with the remaining comb, which is again thermal~\cite[Def.~8, Thm.~6]{Milz26}. Causal order is thus derived from equilibrium preservation. Theorem~\ref{ar:thm:free} starts from causally ordered processes and uses thermal operations instead: its condition~(ii) is a dilation-level, approximate form of their Gibbs-insertion condition, and in condition~(iii) ``thermal'' means energy-conserving unitaries on one retained Gibbs bath, not Gibbs preservation. (An $n$-slot process here is an $(n-1)$-slot supermap in their counting.) The processes that are free in our sense are strictly fewer than theirs; this already follows from the single-channel case~\cite{FOR15,HaagerupMusat11,DingFengHu21}. Milz et al.\ suggest investigating analogous constructions for more restrictive frameworks such as thermal operations, and ask whether causal order can again be derived~\cite[Sec.~VI]{Milz26}. Theorem~\ref{ar:thm:free} gives the thermal-operation counterpart of their thermality condition for processes whose causal order is given; it does not derive causal order.

\begin{proposition}[Order costs work]\label{ar:prop:order}
Let $\Pp_{p,V}$ read a bit $a$ from a qubit $A$ and release it, and then apply $V^{a\oplus r}$ to a second qubit $Q$, where $r$ is a fresh coin equal to $1$ with probability $p$ and $V\in\{X,Z\}$; the Hamiltonians are diagonal, and for $V=X$ the qubit $Q$ has a degenerate one. As one channel on $AQ$, $\Pp_{p,V}$ is a thermal operation. If $a$ must be released before $Q$ arrives, its real-time work rate is $\kB T\ln2\,[1-h(p)]$ per episode; for $V=X$ the bound is already forced by a classical observer.
\end{proposition}

The cheapest device keeps a masked record. Before the episode it prepares a bit $M$ equal to $1$ with probability $p$; at the first slot it adds $a$ into $M$ and releases $a$; at the second it applies $V^M$ to $Q$. The record $M=a\oplus r$ shares $1-h(p)$ bits with the released $a$, and that correlation must be erased after $a$ has left. A device that held $a$ until $Q$ arrived could uncompute the record for free.

Beyond time-covariant processes, the law holds for every finite process when the controller need only return approximately (Theorem~\ref{ar:thm:noncov}): on a rare branch the device records the observer's whole strategy, which tells it how much surplus energy to release as heat without touching any released output. Delayed channels and a correlated process with hidden memory are worked examples (Section~\ref{ar:sec:examples}).

\paragraph{Where the price comes from.} A device that answers in real time may have to keep, in its memory and its baths, correlations with outputs it has already released. It cannot erase those correlations against outputs that are gone. For classical systems the same tension appears in the link between nonpredictive memory and dissipation~\cite{Still12} and in the finding that agents that extract work efficiently must balance prediction against forgetting~\cite{Fiderer26}. A classical channel needs no such correlation on inputs diagonal in its basis, and coherent inputs do not raise its optimized cost, so it pays nothing extra (Appendix~\ref{app:channels}). A quantum channel may need some, and that is the surcharge of Appendix~\ref{app:channels}. A process with memory may have it forced on it by the task itself, classical or not, as is known for classical information ratchets, generators and agents~\cite{BMC18,LoomisCrutchfield20,Thompson25,Fiderer26}: in Proposition~\ref{ar:prop:order} the device must remember $a$ after releasing it.

\paragraph{Prior work.} Faist, Berta and Brand\~ao implement channels that do not respect time translation with thermal operations and a small source of coherence when the input is a fixed i.i.d.\ state~\cite[Sec.~8.2]{FBB21}, and ask whether thermal operations can implement every i.i.d.\ channel universally, which they expect to require an arbitrary amount of coherence~\cite[Sec.~9]{FBB21}. Theorem~\ref{ar:thm:noncov} settles the analogous question for real-time implementation of every finite process, with an unbounded phase reference, in line with their expectation for batch; their batch question, for a clock that must return, is answered in Appendix~\ref{app:batch}. Lie, Son, Boes, Ng and Wilming show that a channel that fixes the Gibbs state is a thermal operation exactly when it has a dilation, by an arbitrary unitary and environment state, that leaves the environment unchanged when the input is the Gibbs state; energy conservation, a Gibbs bath and time covariance then follow~\cite[Def.~1, Lemma~3, Prop.~4]{LieEtAl26}. Theorem~\ref{ar:thm:free} is a robust form of their criterion for processes with several steps: a small defect at Gibbs inputs implies closeness to a thermal implementation, with constants that do not depend on the bath. Theorems~\ref{ar:thm:law} and~\ref{ar:thm:free} give the real-time cost of every finite process whose inputs and outputs carry the same Hamiltonians, and characterize the processes whose cost is zero. For Markovian dynamics, Luo et al.\ show that thermal Lindbladians are exactly those realized by energy-conserving collision models with fresh thermal units~\cite{LuoTrinity26}; the free processes of Theorem~\ref{ar:thm:free} use one bath retained across all slots. For repeated independent uses of one channel at zero Hamiltonian, the rate of Theorem~\ref{ar:thm:law} is the purity cost of~\cite{Douglas26closed}, the zero set of Theorem~\ref{ar:thm:free} is the closure of the finitely factorizable channels~\cite{Memory}, and a strong converse for purity is in~\cite{Flagship}; one channel at every temperature is treated in Appendix~\ref{app:channels}.

\paragraph{Work extraction.} Zambon and Adesso study work extraction from non-Markovian multitime processes~\cite{ZambonAdesso}. An experimenter who stores all outputs until the end can extract more work than one who extracts step by step, by an amount bounded by the non-Markovianity of the process~\cite[Thm.~2]{ZambonAdesso}; this is the extraction-side counterpart of the gap between batch and real-time operation. Since, unlike for states and channels, the work cost of implementing a comb is not known to equal the work it can yield, they take as the cost of a comb the least work value of an ancilla state plus the asymptotic costs~\cite{FBB19,FBB21} of the channels in a dilation. Theorem~\ref{ar:thm:law} concerns a different quantity, the cost of implementing a process in real time, episode after episode, against every observer; it does not settle their question. Fiderer, Barth, Smith and Briegel define the work capacity of a classical environment with memory, the largest rate at which an agent can extract work from it, and ask whether, for channels with memory, it can match the least work the environment must spend to implement its channel~\cite{Fiderer26}. Theorem~\ref{ar:thm:law} gives the least cost of implementing a finite process in real time when it must be accurate against every observer; against a given agent the cost can be lower, and their question remains open. Huang et al.\ define the time-ordered free energy, the most work an agent without persistent quantum memory can extract online from a temporally correlated sequence of quantum states~\cite{HuangTOFE}; that is extraction under a causal constraint, whereas we price implementation with arbitrary quantum memory.

\paragraph{Classical modularity.} For classical processes, the price of operating step by step rather than on the whole process at once is the modularity dissipation of Boyd, Mandal and Crutchfield, the loss of global correlations under local control~\cite{BMC18}, which Loomis and Crutchfield trace to temporal locality~\cite{LoomisCrutchfield20}; its simplest instance is the Landauer--Bennett price of a record whose original has left~\cite{Bennett73,Bennett82}. Proposition~\ref{ar:prop:order} gives its exact rate in a process that is free when its outputs may wait.

\paragraph{Online agents.} Thompson et al.\ prove their classical online cost optimal among causal agents, whose memory is a function of the past~\cite{Thompson25}. A classical agent that commits its randomness in advance lies outside that class and can pay less (Section~\ref{ar:sec:examples}). When the inputs stay on the tape, this is the analogue, for a transducer whose outputs depend on informative inputs, of the known efficiency of retrodictive pattern generators~\cite{BMRC17,BMC18,LoomisCrutchfield20}; when each output overwrites its input, even this agent pays $1-h(p)$ (Section~\ref{ar:sec:examples}).

\paragraph{Outline.} Section~\ref{ar:sec:setting} sets up the model and defines $\Kb$. Section~\ref{ar:sec:law} proves the work law, Section~\ref{ar:sec:converse} the strong converse and Section~\ref{ar:sec:free} Theorem~\ref{ar:thm:free}. Section~\ref{ar:sec:noncov} drops time covariance, and Section~\ref{ar:sec:examples} treats channels, the ordered and delay processes, and classical agents. The longer constructions are in the appendices. Section~\ref{sec:law} places this law among general schedules of arrivals, releases and loans.

% =====================================================================
